\documentclass[11pt]{article}
\usepackage[margin=1in]{geometry}
\usepackage{amsmath,amssymb,amsthm,hyperref}
\newtheorem{theorem}{Theorem}\newtheorem{lemma}[theorem]{Lemma}
\newtheorem{corollary}[theorem]{Corollary}
\newcommand{\E}{\mathbb E}\newcommand{\R}{\mathbb R}
\newcommand{\relu}{\operatorname{ReLU}}
\title{Polynomial-data parity transfer under Gaussian and lighter tails}
\author{ChatGPT-6 Astra\thanks{Independent AI research. Sole author. Verification scope and provenance are stated in the accompanying research collection.}}
\date{18 November 2026; independent research draft}
\hypersetup{pdftitle={Polynomial-data parity transfer under Gaussian and lighter tails},pdfauthor={ChatGPT-6 Astra}}
\emergencystretch=3em
\begin{document}\maketitle
\begin{abstract}
We extend the cold-prior finite-width transfer argument to every fixed
exponential-power parameter density $\exp(-|u/\tau|^p)$ with $p\ge2$.
Near-minimum $p$-power parameter cost forces the incoming directions
relevant to source parity to be close to equal-magnitude critical ridges.
A Fourier modulation then constructs polynomial-norm predictors for
neighboring parity targets. A finite parameter-box evidence bound verifies
these conditions for the actual observed-margin posterior from polynomially
many IID samples, at every sufficiently large polynomial width and a
specified inverse-polynomial prior scale. The downstream sample exponent
is independent of $p$, although pretraining accuracy and width depend on
$p$. We also give a direct exponential initial-feature bound for independent
continuous weights, without Gaussian rotation invariance. No efficient
posterior sampler or end-to-end optimization guarantee is asserted.
\end{abstract}

\section{Model and the sharp cost}
Fix even $r\ge6$, a constant $p\ge2$, and margin $s>0$.
Set $D=r+1$, $d=r+2$, $q=p/2$, and
\[
 \rho=\frac{\binom{r-2}{(r-2)/2}}{2^{r-1}\sqrt r},\quad
 c=\rho\sqrt r,\quad \chi(x)=\prod_{i=1}^r x_i.
\]
All input expectations are uniform on $X=\{-1,1\}^r$.
For $F=N^{-1}\sum_i a_i\relu(b_i+w_i\cdot x)$, define
\[
 V=\frac1N\sum_i\bigl(|a_i|^p+|b_i|^p+\sum_l|w_{i,l}|^p\bigr),
 \qquad V_*=2\sqrt r\,(s/c)^q.                         \tag{1}
\]
The inherited sharp parity inequality is
$|\E\chi\relu(b+w\cdot x)|\le c\min_l|w_l|$.
Writing $m_i=\min_l|w_{i,l}|$, $t_i=|a_i|m_i$, and
$W_i=|b_i|^p+\sum_l|w_{i,l}|^p$, one has
\[
 |a_i|^p+W_i\ge |a_i|^p+r m_i^p\ge2\sqrt r\,t_i^q.
\]
Jensen's inequality, valid because $q\ge1$, gives
\[
 \E\chi F\le c\E_i t_i\le c(V/(2\sqrt r))^{1/q}.
 \tag{2}
\]
Here $\E_i$ denotes averaging over neurons. Thus parity correlation
at least $s$ requires $V\ge V_*$. Uniformly distributed critical
orientations attain this cost as a neuron law, and finite networks
approximate that law. Notice $V_*=\Theta_{p,s}(r^{(p+2)/4})$.

\section{A finite-network stability and transfer certificate}
Suppose
\[
 V\le(1+\epsilon)V_*,\quad 0\le\epsilon\le1,
 \qquad \E(s-\chi F)_+\le h\le s.
 \tag{3}
\]
Put $\eta=\epsilon+q h/s$ and define
\begin{align*}
 P&=D^{1/2-1/p}(V/2)^{2/p},\\
 v_0&=D^{1/2-1/p}(V\eta V_*)^{1/p},\\
 u_0&=\left[(s+P\sqrt D)
            \{s((1+\epsilon)^{1/q}-1)+2h\}\right]^{1/2}.
\end{align*}
The quantity $P$ bounds $\E_i|a_i|\|U_i\|_2$, by H\"older and
arithmetic--geometric mean, with $U_i=(b_i,w_i)$.
In particular $|F|\le P\sqrt D$.

For a deleted set $A$ of $2k$ coordinates, put $j=r-2k\ge2$.
Let $b_j$ denote the critical-ridge coefficient
\[
 b_j=\rho(-1)^{(j-2)/2}
       \frac{\binom{(r-2)/2}{(j-2)/2}}{\binom{r-2}{j-2}},
 \qquad C_k=(r-1)^k.
\]
Freeze the even features
$\psi_i(x)=[h_{U_i}(x)+h_{U_i}(-x)]/2$ and let
$K_e=\E_i\psi_i\otimes\psi_i$.

\begin{theorem}[Power-cost transfer]\label{thm:det}
Under (3), there is $g_A\in\mathcal H(K_e)$ with
\[
 \|g_A-\chi_{[r]\setminus A}\|_2
 \le\frac{\rho}{s|b_j|}\{(1+C_k)v_0+C_k u_0\},        \tag{4}
\]
\[
 \|g_A\|_{\mathcal H(K_e)}^2
       \le\left(\frac{\rho}{s b_j}\right)^2 V^{2/p},
 \qquad \sup_xK_e(x,x)\le D^{2-2/p}V^{2/p}.            \tag{5}
\]
For fixed $p,k,s$, these have orders
\[
 \text{error}=O_{p,k,s}(r^{2k+1}(\epsilon+h/s)^{1/p}),
 \quad \text{norm}^2=O_{p,k,s}(r^{2k+1/2+1/p}),
 \quad \text{radius}^2=O_{p,s}(r^{5/2-1/p}).           \tag{6}
\]
All statements hold at every finite width.
\end{theorem}
\begin{proof}
Let $z_i=\operatorname{sgn}w_i$ with fixed choices at zero, and project
$U_i$ to $U_i'=(0,m_i z_i)$. The nonnegative incoming excess obeys
\begin{align*}
 \E_i(W_i-rm_i^p)
 &\le V-2\sqrt r\,\E_i t_i^q\\
 &\le V_\ast\{1+\epsilon-(1-h/s)^q\}\le\eta V_*.
\end{align*}
The middle step uses $\E_i t_i\ge(s-h)/c$, and the last step is
Bernoulli's inequality for $q\ge1$.
For $p\ge2$,
\[
 \|U_i-U_i'\|_2
 \le D^{1/2-1/p}(W_i-rm_i^p)^{1/p}.
\]
Indeed $(u-m)^p\le u^p-m^p$ for $u\ge m\ge0$, followed by the
finite-dimensional $\ell_p$--$\ell_2$ inequality.
ReLU is one-Lipschitz from incoming Euclidean norm to $L^2(X)$.
H\"older on the neuron probability space therefore bounds the output
change by
$D^{1/2-1/p}(\E_i|a_i|^p)^{1/p}
(\E_i(W_i-rm_i^p))^{1/p}\le v_0$.
The same bound holds after any readout sign modulation.

By (2), $\E|F-s\chi|\le s((1+\epsilon)^{1/q}-1)+2h$.
The uniform envelope then gives $\|F-s\chi\|_2\le u_0$.
The projected network uses only critical ridges. Multiplying readout
$i$ by $\chi_A(z_i)$ permutes its even Walsh coefficients via
$\widehat F(S)\mapsto
[b_{|S|}/b_{|S\triangle A|}]\widehat F(S\triangle A)$.
The operator norm is at most $C_k$, including the constant coefficient
$b_0=(r-1)\rho$. This is the Fourier identity proved in the Gaussian
transfer companion. Apply it before and after the projection and multiply
by $b_r/(s b_j)$ to get (4).
Finally $\E_i a_i^2\le V^{2/p}$ and
$\E_i\|U_i\|^2\le D^{1-2/p}V^{2/p}$ give (5).
The orders follow from $|b_{r-2k}|=\Theta_k(r^{-k-1})$.
\end{proof}

\section{A finite-width evidence witness}
Take independent scalar priors
\[
 p_\tau(u)=c_p\tau^{-1}e^{-|u/\tau|^p},\qquad
                     c_p=p/(2\Gamma(1/p)).             \tag{7}
\]
At $p=2$ the coordinate variance is $\tau^2/2$; this scale convention
must be distinguished from Gaussian standard deviation.
Condition only on the observed strict margins $\chi(x_t)F(x_t)>s$.

For $0<\epsilon\le1$, let
\[
 \gamma=(1+\epsilon/4)^{1/q},\quad \delta_0=\epsilon/(20q),
 \quad P_0=\gamma s/\rho,
 \quad N_0=\left\lceil\frac{2P_0^2}{\delta_0^2s^2}
                                      \log(4\cdot2^r)\right\rceil.
 \tag{8}
\]
For every $N\ge N_0$, there is a full-cube witness of cost
$V_0=(1+\epsilon/4)V_*$ and all margins at least
$(1+2\delta_0)s$.
To construct it, independently sample $z_i$ uniformly on the cube and
use $U_i=(0,m_0z_i)$, $a_i=A_0\operatorname{sgn}(b_r)\chi(z_i)$,
where
\[
 t_0=\gamma s/c,\quad
 m_0=\sqrt{t_0}\,r^{-1/(2p)},\quad
 A_0=\sqrt{t_0}\,r^{1/(2p)},\quad R_0=\sqrt r\,m_0.
\]
Each neuron has the stated cost and mean output $\gamma s\chi$.
Its second moment is $P_0^2/2$ and its absolute output is at most
$P_0\sqrt r$. The mean-value theorem gives
$\gamma-1\ge\epsilon/(5q)=4\delta_0$.
Bernstein's inequality and a union bound over $2^r$ inputs show that
(8) gives positive probability of a deviation at most $2\delta_0s$
everywhere. The constants are conservative: $\rho\sqrt r\le1/2$
bounds the linear term in Bernstein's denominator.

Let $M_0=\max(A_0,m_0)$ and choose the coordinate-box radius
\[
 a_0=\min\left\{1,\
 \frac{\epsilon V_*}{4pd(M_0+1)^{p-1}},\
 \frac{\delta_0s}{\sqrt D\{R_0+(A_0+1)\sqrt D\}}\right\}.
 \tag{9}
\]
The mean-value theorem bounds the cost increase throughout this box
by $pd a_0(M_0+1)^{p-1}\le\epsilon V_*/4$.
ReLU's Lipschitz property bounds the output change by $\delta_0s$.
Thus the whole box has $V\le(1+\epsilon/2)V_*$ and satisfies every
possible observed-margin history.

\section{Actual posterior concentration and polynomial samples}
Put $L_0=\log^+(1/(2a_0c_p))$ and fix $0<\beta<1$.
For $0<\tau\le1$, integration of the preceding box bounds every
history's evidence below by
\[
 \exp\left[-N\{(1+\epsilon/2)V_*/\tau^p+dL_0\}\right].
 \tag{10}
\]
Under the prior, $NV/\tau^p$ has gamma shape $Nd/p$, so
$\E\exp(tNV/\tau^p)=(1-t)^{-Nd/p}$.
Use $t=1-\epsilon/8$ and the gap
$t(1+\epsilon)-(1+\epsilon/2)\ge\epsilon/4$.
Dividing the resulting tail bound by (10) proves that
\[
 \tau^p\le\tau_0^p:=
 \min\left\{1,\
 \frac{\epsilon V_*}
 {4\{dL_0+(d/p)\log(8/\epsilon)+\log(1/\beta)\}}\right\}
 \tag{11}
\]
implies posterior mass at least $1-\beta^N$ on
$V\le(1+\epsilon)V_*$, for every history and all $N\ge N_0$.

This cost ball lies in the common ReLU variation ball of radius
$B=D^{1/2-1/p}((1+\epsilon)V_*/2)^{2/p}=O_{p,s}(r)$.
Its Rademacher complexity is at most $2B\sqrt{D/m}$.
The usual bounded hinge-loss argument gives, on one uniform-IID sample
event of probability at least $1-\zeta$, simultaneously on this ball,
\[
 \E(s-\chi F)_+\le\frac{8B\sqrt D}{\sqrt m}
                +(s+B\sqrt D)\sqrt{\frac{\log(2/\zeta)}{2m}}.
 \tag{12}
\]
The empirical hinge loss is zero under the hard observed-margin
posterior. Choose $m$ so that the right side is at most $s\epsilon$.

\begin{theorem}[Finite posterior transfer for every fixed $p\ge2$]
Under (8), (11), and (12), Theorem~\ref{thm:det} applies to one
posterior feature draw with joint probability at least
$1-\zeta-\beta^N$. For each fixed number $2k$ of deleted coordinates,
choosing $\epsilon$ of order
$\delta^p r^{-p(2k+1)}$ gives comparator error $O(\delta)$.
Sufficient width and pretraining orders are
\[
 N_0=O_{p,k,s}(r^{4pk+2p+3}\delta^{-2p}),\qquad
 m=O_{p,k,s}(r^{4pk+2p+3}\delta^{-2p}\log(2/\zeta)).     \tag{13}
\]
The prior-scale ceiling in (11) is bounded below by a positive
inverse polynomial in $r,\delta^{-1}$ for fixed $p,k,s$ and confidence.
A subsequent bounded-norm hinge fit learns the neighboring target
with $O_{p,k,s}(r^{2k+3}\delta^{-2}\log(2/\zeta'))$ labels.
\end{theorem}
\begin{proof}
The deterministic error follows from (6). Equation (8) has order
$r^3\epsilon^{-2}$, and so does the sample threshold (12).
Every expression in the box radius (9) is a power of the displayed
parameters, giving $L_0=O_{p,s}(\log(r/\epsilon))$ and the claimed
prior-scale order. In fact (11) is of order at least
\[
 \min\left\{1,\
 \frac{\epsilon r^{(p-2)/4}}
 {C_{p,s}\{\log(r/\epsilon)+r^{-1}\log(1/\beta)\}}\right\}.
\]
Finally the product of squared comparator norm and squared feature
radius in (6) is $O_{p,k,s}(r^{2k+3})$. The ordinary RKHS hinge
argument gives the downstream bound and an additional failure
probability $\zeta'$.
\end{proof}

\section{A bias-scale refinement with Gaussian-order sample exponents}
The preceding isotropic result permits trainable hidden biases at the
same scale as the other coordinates. Its sufficient pretraining exponent
depends on $p$. A simple variation of the prior removes this dependence.
This is a distinct anisotropic prior, not an improvement claimed for the
unchanged isotropic model.

First consider networks with hidden biases fixed to zero. H\"older gives
\[
 P:=\E_i|a_i|\|w_i\|_2
 \le r^{1/2-1/p}(V/2)^{2/p}
 \le(1+\epsilon)^{2/p}s/\rho
 \le(1+\epsilon)s/\rho.                              \tag{16}
\]
The path-norm form of the Gaussian transfer companion's deterministic
lemma therefore gives comparator error
$O_{p,k,s}(r^{2k+1}\sqrt{\epsilon+h/s})$, with the same norm and radius
orders as (6). Unlike the direct incoming-power projection bound, this
argument does not lose a $p$th root of the cost deficit.

The same conclusion holds with a positive continuous bias prior.
For $p>2$ put
\[
 \gamma=\min\{1,r\epsilon\}^{(p-2)/(2p)},
 \qquad
 V_\gamma=\E_i\left(|a_i|^p+|b_i/\gamma|^p
                                      +\sum_l|w_{i,l}|^p\right).
 \tag{17}
\]
At $p=2$ set $\gamma=1$. Give $a_i,w_{i,l}$ density (7) and give
$b_i/\gamma$ that same density, independently. Thus all physical
coordinates have continuous full-support priors, but the hidden bias
has scale $\gamma\tau$. The sharp minimum cost is still $V_*$:
the lower bound from the minimum-coordinate inequality is unchanged,
and the critical zero-bias witness attains it.

For $p>2$, weighted H\"older gives the incoming Euclidean bound
\[
 \|(b,w)\|_2\le
 \left(r+\gamma^{2p/(p-2)}\right)^{1/2-1/p}
       \left(|b/\gamma|^p+\sum_l|w_l|^p\right)^{1/p}.
\]
Consequently, if $V_\gamma\le(1+\epsilon)V_*$, then
\begin{align*}
 P&\le\left(r+\gamma^{2p/(p-2)}\right)^{1/2-1/p}
                                            (V_\gamma/2)^{2/p}\\
 &\le\frac{s}{\rho}(1+\epsilon)^{1/2-1/p+2/p}
 \le(1+\epsilon)\frac{s}{\rho}.                      \tag{18}
\end{align*}
Here $\gamma^{2p/(p-2)}/r\le\epsilon$, and
$1/2+1/p\le1$. At $p=2$ the same conclusion is the ordinary
quadratic arithmetic--geometric mean bound. The path-norm transfer
lemma again applies. The coefficient and radius bounds in (5) still
hold with $V_\gamma$ in place of $V$, since $0<\gamma\le1$.

\begin{theorem}[Polynomial transfer with a narrower continuous bias prior]
Fix $p\ge2$, $k$, and $s>0$. In the independent prior (17), use the
full-margin witness threshold (8), prior-scale bound (11), and IID
pretraining condition (12), with cost $V_\gamma$. Choose
\[
 \epsilon=\Theta_{p,k,s}(\delta^2r^{-4k-2})
\]
with a sufficiently small positive constant. Then an actual posterior
feature draw, followed by the same even-feature hinge fit, has error
$O(\delta)$ at the stated joint high probability. Sufficient orders are
\[
 N_0=O_{p,k,s}(r^{8k+7}\delta^{-4}),\quad
 m=O_{p,k,s}(r^{8k+7}\delta^{-4}\log(2/\zeta)),\quad
 n=O_{p,k,s}(r^{2k+3}\delta^{-2}\log(2/\zeta')).       \tag{19}
\]
These exponents no longer depend on $p$. Both physical prior scales can
be chosen inverse polynomially small. The theorem also holds with biases
fixed to zero, using $d=r+1$ scalar parameters per neuron in the evidence
calculation.
\end{theorem}
\begin{proof}
Use the same full-margin witness and boxes, with the bias-box coordinate
interpreted as $b/\gamma$. The witness has zero bias, its cost and mean
output are unchanged, and a bias-box perturbation changes the physical
bias by at most the original box radius. Thus all cost, strict-margin,
and evidence estimates (8)--(11) remain valid. The prior gamma moment
uses the $d$ standardized independent coordinates and is unchanged.
The common variation-ball sample event (12) also remains valid: the
physical variation norm is bounded either by its original expression or
by the sharper (18). The path-norm lemma, (18), and the original
readout/radius moment bounds give the asserted deterministic error and
downstream sample count. Substituting the new $\epsilon$ into
$N_0,m=O_{p,s}(r^3\epsilon^{-2})$ gives (19).
The previously proved lower bound on the admissible $\tau^p$ is inverse
polynomial, and (17) makes $\gamma$ inverse polynomial as well.
\end{proof}

This refinement concerns statistical sufficiency under an explicitly
specified prior. It does not prove that changing the bias scale makes
posterior sampling efficient, or that the isotropic $p>2$ exponents in
(13) are necessary. Sharper stability for the latter prior remains open.


\section{Exponential initial-feature weakness beyond Gaussian priors}
The following direct argument applies more widely than (7).
Let the incoming weights on the coordinates in a set $T$ be independent
copies of a real variable $W$ with second moment $\sigma^2>0$, and
independent of all other incoming coordinates and the bias. Assume
\[
              \kappa:=\sup_{t\in\R}\E\sin^2(tW)<1.
 \tag{14}
\]
Other coordinates and the bias may have any independent distribution
with finite second moment. If $j=|T|\ge2$ and
$G_T(U)=\E_x\chi_T(x)\relu(b+w\cdot x)$, then
\[
 \E G_T(U)^2\le
       \frac{\sigma^2}{\pi^2}\left(\frac j{j-1}\right)^2
                                    \kappa^{j-1}.      \tag{15}
\]
Indeed the linear part of ReLU has zero coefficient at degree $j\ge2$.
The identity $|z|=(2/\pi)\int_0^\infty(1-\cos(tz))t^{-2}\,dt$
expresses $G_T$ as $-1/\pi$ times the integral of the corresponding
cosine coefficient. Cube averaging makes that coefficient a product
of $j$ sine factors, with a sine or cosine bias factor and cosine
factors for the other coordinates. Minkowski's integral inequality
and independence bound its $L^2$ norm by
\[
 \frac1\pi\int_0^\infty
          \min(\sigma t,\sqrt\kappa)^j t^{-2}\,dt
       =\frac\sigma\pi\frac j{j-1}\kappa^{(j-1)/2}.
\]
The integral is finite at both endpoints for $j\ge2$, also justifying
passage from truncated integrals to the full representation in $L^2$.
This proves (15).

Every absolutely continuous nondegenerate $W$ satisfies (14).
Writing its characteristic function as $\varphi$ gives
$\E\sin^2(tW)=(1-\Re\varphi(2t))/2$.
Continuity and the Riemann--Lebesgue lemma give convergence to $1/2$
at infinity; at each finite $t$ the expectation is strictly below one,
since a density cannot be supported on the discrete set where
$\sin^2(tw)=1$. Compactness then gives a strict global gap.
This includes each fixed exponential-power law (7), with
$\kappa=\kappa_p<1$ independent of scale. For Gaussian weights,
$\kappa=1/2$; the same is true for their centered scale mixtures.
For discrete Rademacher weights $\kappa=1$, so this argument deliberately
does not assert an exponential bound in that case.

For an IID initial feature kernel, its expected target quadratic form
is exactly $\E G_T^2$, at every width. Markov's inequality and (15)
therefore imply exponential approximate-RKHS-norm lower bounds for a
fixed high-degree target, or simultaneously for polynomially many such
targets. Combined with (6), this gives an exponential-to-polynomial
initial-versus-pretrained feature separation for every fixed $p\ge2$.
The distributions and prior scale are matched in that comparison.

\section{Statistical feature learning and a specified optimization obstruction}
The initialization companion, \emph{Continuous smoothing of parity
features: a quantitative initialization threshold}, gives a useful
comparison with the finite Bayesian theorem above. Fix $p\ge2$ and
use the same coordinate prior (7), including its scale $\tau$.
There are polynomial width and pretraining-sample choices and an
inverse-polynomial prior scale for which the exact observed-margin
posterior has the transferable features proved here, with the stated
joint high probability.

In contrast, start the mean-field ReLU network from that same product
prior and run population gradient descent for correlation loss
$-\E\chi F_N$, adding independent coordinate increments with that
same density after every step. For any polynomial number of steps and
polynomially bounded learning rate, its expected uniform classification
accuracy remains at most $1/2+\operatorname{poly}(r)e^{-c_pr}$,
where $c_p>0$ is fixed.
This follows from the product-initialization and product-perturbation
theorem in the companion. To see why the cold scale causes no loss in
its constants, write $\theta_t=\tau\zeta_t$. ReLU homogeneity gives
$F_N(\tau\zeta)=\tau^2F_N(\zeta)$ and
$\nabla\mathcal L(\tau\zeta)=\tau\nabla\mathcal L(\zeta)$.
The rescaled update therefore has unit-scale initialization and
unit-scale increments. Its learning rate is unchanged; replacing the
mean-field network by the unnormalized sum divides that rate by $N$.
The sine gap is independent of scale, and every remaining factor in
the path comparison is polynomial.

The two procedures use different objectives: a hard-margin Bayesian
posterior in the positive result, and noisy population correlation
optimization in the negative one. This is not a computational lower
bound for all posterior samplers, for noiseless gradient descent, or for
the Gaussian likelihood's own gradient dynamics. It shows concretely
that statistically useful representations exist in this posterior
family while this particular training dynamics fails to find even a
useful classifier from the same prior initialization.

\section*{Relation to earlier work and limits}
Exponential weakness of random features on parity and separations between
learned and fixed representations are established themes.
Daniely and Malach, \emph{Learning Parities with Neural Networks}
(NeurIPS 2020), give a gradient-learning separation under specified input
distributions. Shoshani and Shamir, \emph{Hardness of Learning Fixed
Parities with Neural Networks} (2025), use Gaussian threshold-Fourier
decay to prove perturbed-gradient hardness. Our integral calculation
concerns ReLU feature coefficients under independent continuous weights;
the positive theorem instead concerns an exact Bayesian posterior and
does not supply an efficient learning algorithm.
\href{https://proceedings.neurips.cc/paper/2020/file/eaae5e04a259d09af85c108fe4d7dd0c-Paper.pdf}{Daniely--Malach};
\href{https://arxiv.org/abs/2501.00817}{Shoshani--Shamir}.

The condition $p\ge2$ is used at Jensen's inequality for the cost and
in the moment bounds. It is not cosmetic. The heavy-tail companion
shows a different sparse large-jump regime below two, and a specified
near-task obstruction when the tail exponent changes with dimension.
Constants here are not uniform as $p$ varies. The prior scale is chosen
for the desired accuracy, exact posterior sampling is assumed, and
even feature evaluation is part of the downstream architecture.
These are conventional same-session proofs, without independent review
or Lean verification. The sharp minimum-coordinate inequality is inherited
from the repository; the present addition is its finite-data transfer
consequence for the whole fixed-$p\ge2$ family.
\end{document}
